\documentclass[11pt,a4paper]{report}

\usepackage[utf8]{inputenc}
\usepackage[T1]{fontenc}
\usepackage{lmodern}
\usepackage{amsmath,amssymb,amsfonts}
\usepackage{graphicx}
\usepackage{booktabs}
\usepackage{hyperref}
\usepackage{geometry}
\usepackage{enumitem}

\geometry{margin=1in}
\hypersetup{
    colorlinks=true,
    linkcolor=blue,
    citecolor=blue,
    urlcolor=blue
}

\title{Local and Global Equilibrium Effects\\in Discrete Graph Attention Diffusion (DGAD)}
\author{DGAD Discussion Notes}
\date{June 2026}

\begin{document}

\maketitle
\tableofcontents
\newpage

%==============================================================================
\chapter{Introduction}
%==============================================================================

Discrete Graph Attention Diffusion (DGAD) iteratively smooths node features on a graph by moving each node toward an attention-weighted mixture of its neighbors. A single diffusion step in \texttt{GNDLayer} takes the form
\begin{equation}
    x_i^{(k+1)} = (1-\tau)\, x_i^{(k)} + \tau \sum_{j \in \mathcal{N}(i)} \alpha_{ij}^{(k)}\, x_j^{(k)},
    \label{eq:node-tau}
\end{equation}
where $\tau$ is the scalar \texttt{time\_increment} (default $0.2$), and $\alpha_{ij}$ are neighborhood-normalized attention weights. Stacking $K$ such steps (\texttt{num\_steps}) yields final node embeddings suitable for classification, clustering, and other downstream tasks.

A recurring phenomenon in this process is the coexistence of \textbf{local} and \textbf{global equilibrium effects} in node embeddings. This document collects two complementary interpretations of that phenomenon, discusses how to exploit the trade-off for community separation, surveys methods for designing per-edge step sizes $\tau_{ij}$, and finally proposes a KNN-specific construction based on reverse-degree asymmetry.

%==============================================================================
\chapter{Equilibrium Effects in DGAD}
%==============================================================================

\section{The Core Mechanism}

Each diffusion step balances two forces:
\begin{itemize}[noitemsep]
    \item \textbf{Retention}: the term $(1-\tau)\, x_i^{(k)}$ preserves the node's current identity;
    \item \textbf{Mixing}: the term $\tau \sum_j \alpha_{ij} x_j^{(k)}$ pulls the node toward its attention-weighted neighborhood aggregate.
\end{itemize}

At a \emph{fixed point} (equilibrium) for node $i$, assuming $\tau > 0$,
\begin{equation}
    x_i = \sum_{j \in \mathcal{N}(i)} \alpha_{ij}\, x_j.
    \label{eq:local-fixed-point}
\end{equation}
The embedding equals the attention-weighted average of its neighbors. How one labels this ``local'' versus ``global'' depends on the scale at which equilibrium is observed.

\section{Interpretation I: Topology and Hop Count}

\subsection{Local Equilibrium (One-Hop Scale)}

\textbf{Local equilibrium} refers to the per-node, per-step attractor defined by Equation~\eqref{eq:local-fixed-point}. Within a single step, each node moves toward consensus with its immediate neighbors only. The step size $\tau$ controls how far the node travels in one update:
\begin{itemize}[noitemsep]
    \item small $\tau$: embeddings retain individual structure longer;
    \item large $\tau$: faster convergence to the one-hop neighbor mixture.
\end{itemize}

Attention determines \emph{which} neighbors matter and \emph{how much}; it shapes the local attractor itself rather than applying uniform averaging.

\subsection{Global Equilibrium (Multi-Hop Scale)}

\textbf{Global equilibrium} emerges after $K$ repeated local updates. Although each step mixes only one-hop neighbors, information propagates across the graph over multiple steps---analogous to the continuous-time dynamics
\begin{equation}
    \frac{dx}{dt} = \mathrm{Agg}(x) - x,
\end{equation}
whose long-time behavior drives connected components toward graph-wide harmonization (the familiar oversmoothing limit).

\subsection{Correspondence Between Scales}

\begin{table}[h]
\centering
\caption{Two temporal--spatial scales of equilibrium.}
\begin{tabular}{@{}lll@{}}
\toprule
\textbf{Aspect} & \textbf{Local} & \textbf{Global} \\
\midrule
Spatial scope   & 1-hop neighborhood        & Multi-hop / full graph \\
Temporal scope  & Single step               & Accumulated over $K$ steps \\
Manifestation   & Node individuality vs.\ neighbors & Distant nodes influence one another \\
Control knobs   & $\tau$ (per-step mixing)  & $K$ (\texttt{num\_steps}) \\
\bottomrule
\end{tabular}
\end{table}

Default settings ($\tau = 0.2$, $K = 8$) strike a practical balance: enough local smoothing without immediate global collapse.

\section{Interpretation II: Community Structure}

\subsection{Local Equilibrium Within Communities}

On graphs with community structure, an alternative and often more task-relevant reading applies:
\begin{itemize}[noitemsep]
    \item \textbf{Local equilibrium} = equilibrium \emph{within} a community $C$.
\end{itemize}

Inside $C$, edges are dense, features are homophilous, and attention softmax is dominated by intra-community neighbors. After a few diffusion steps, nodes in $C$ converge toward a \textbf{community-specific consensus}---similar embeddings inside $C$, while distinct communities remain separable. This occurs on a \emph{short} diffusion time scale because information circulates primarily within $C$.

\subsection{Global Equilibrium Between Communities}

\begin{itemize}[noitemsep]
    \item \textbf{Global equilibrium} = equilibrium \emph{between} communities.
\end{itemize}

Inter-community edges are sparse; bridge and boundary nodes mediate cross-community flow. Cross-community neighbors receive lower attention mass, so inter-community information exchange is slow and requires more steps. Global equilibrium means community-level representations settle into a coordinated graph-wide configuration---not necessarily identical vectors, but stable relative positions across communities.

\subsection{Why DGAD Exhibits Both Scales}

The update in Equation~\eqref{eq:node-tau} naturally splits edges into two regimes:

\begin{table}[h]
\centering
\caption{Edge regimes under community structure.}
\begin{tabular}{@{}lll@{}}
\toprule
\textbf{Connection} & \textbf{Attention pattern} & \textbf{Equilibration speed} \\
\midrule
Intra-community $(i,j \in C)$ & Large weights, many competing neighbors & Fast $\to$ \textbf{local equilibrium} \\
Inter-community $(i \in C,\, j \in C')$ & Small weights, few channels & Slow $\to$ \textbf{global equilibrium} \\
\bottomrule
\end{tabular}
\end{table}

Diffusion is therefore \textbf{multi-scale}:
\begin{quote}
\emph{Few steps} $\Rightarrow$ communities equilibrate internally; inter-community coupling remains weak.\\
\emph{Many steps} $\Rightarrow$ intra-community states are stable; cross-community alignment accelerates.
\end{quote}

\subsection{Unifying the Two Interpretations}

On community-structured graphs, the hop-count and community views largely coincide: one-hop mixing predominantly occurs inside communities, while multi-hop propagation is required for sustained inter-community coupling. The community interpretation explains \emph{why} two time scales appear---because the graph itself is heterogeneous---not merely because step count is large.

%==============================================================================
\chapter{Promoting Local Equilibrium, Slowing Global Equilibrium}
%==============================================================================

\section{Principle for Community Separation}

To better \textbf{distinguish communities}, one should:
\begin{itemize}[noitemsep]
    \item \textbf{Promote local equilibrium}: encourage fast intra-community consensus (tight, coherent clusters);
    \item \textbf{Slow global equilibrium}: limit cross-community homogenization (preserve boundaries).
\end{itemize}

The desirable operating point is often: \emph{intra-community equilibrium has been reached, but inter-community equilibrium has not}.

\section{Hyperparameter Guidance}

Equation~\eqref{eq:node-tau} offers two primary global knobs:

\begin{table}[h]
\centering
\caption{Global controls for the local/global trade-off.}
\begin{tabular}{@{}lll@{}}
\toprule
\textbf{Goal} & \textbf{Direction} & \textbf{Parameter} \\
\midrule
Promote local equilibrium   & Strengthen intra-community mixing (moderately) & $\tau$ (\texttt{time\_increment}) \\
Slow global equilibrium     & Limit cross-community propagation            & $K$ (\texttt{num\_steps}), or reduce $\tau$ \\
\bottomrule
\end{tabular}
\end{table}

\paragraph{Caveat on $\tau$.}
Very large $\tau$ accelerates \emph{both} local and global equilibration and can cause oversmoothing. A robust strategy is:
\begin{enumerate}[noitemsep]
    \item \textbf{Control $K$ first}---the most direct lever on global reach;
    \item \textbf{Tune $\tau$ moderately} in combination with $K$;
    \item \textbf{Prefer fixed topology} (\texttt{edge\_rewire=False}) when the input graph already encodes sparse cross-community links;
    \item \textbf{Use distance-based attention} (\texttt{attention\_type="dist"}) to up-weight similar neighbors and down-weight dissimilar ones, biasing toward ``fast local, slow global'' behavior when inter-community feature gaps are large.
\end{enumerate}

\section{Practical Regimes}

\begin{table}[h]
\centering
\caption{Qualitative regimes.}
\begin{tabular}{@{}ll@{}}
\toprule
\textbf{Setting} & \textbf{Outcome} \\
\midrule
Small $K$ + moderate $\tau$ & Local equilibrium inside communities; global equilibrium delayed --- \textbf{favorable for community resolution} \\
Large $K$ + large $\tau$       & Both equilibria advance; oversmoothing --- \textbf{communities hard to separate} \\
\bottomrule
\end{tabular}
\end{table}

\section{Downstream Task Implications}

\begin{itemize}[noitemsep]
    \item \textbf{Unsupervised clustering}: the ``local yes, global not yet'' window often aligns with better cluster structure.
    \item \textbf{Node classification}: intra-community smoothing supports label consistency; preserving inter-community differences retains class boundaries.
    \item \textbf{Over-diffusion}: performance can degrade not from insufficient smoothing, but from \emph{excessive global equilibrium} that erases discriminative structure.
\end{itemize}

%==============================================================================
\chapter{Per-Edge Step Sizes \texorpdfstring{$\tau_{ij}$}{tau\_ij}}
%==============================================================================

\section{From Scalar to Per-Edge \texorpdfstring{$\tau$}{tau}}

The README describes a single adjustable step size $\tau$ per node update. A natural extension assigns a distinct $\tau_{ij}$ to each directed edge $(j \to i)$, yielding the edge-wise incremental form
\begin{equation}
    x_i^{(k+1)} = x_i^{(k)} + \sum_{j \in \mathcal{N}(i)} \alpha_{ij}\, \tau_{ij}\, \bigl(x_j^{(k)} - x_i^{(k)}\bigr).
    \label{eq:edge-tau}
\end{equation}

Interpretation:
\begin{itemize}[noitemsep]
    \item large $\tau_{ij}$: node $i$ mixes rapidly along edge $(j \to i)$;
    \item small $\tau_{ij}$: that edge contributes little cross-node drift.
\end{itemize}

Constraint: $\tau_{ij} \in [0,1]$. Optionally enforce a per-node budget $\sum_j \alpha_{ij}\tau_{ij} \le \tau_{\max}$ for stability.

\paragraph{Implementation note.}
In the current \texttt{GNDLayer}, \texttt{time\_increment} is applied \emph{after} neighborhood aggregation at the node level. True per-edge $\tau_{ij}$ should enter the weighted aggregation, e.g.,
\begin{verbatim}
weighted = nodes_features_source * attentions_per_edge * tau_per_edge
\end{verbatim}
with node-level propagate optionally reduced to a pure residual or $\tau_i \equiv 1$.

\section{Design Methods}

\subsection{Community Priors (Labels or Detected Communities)}

When communities $c_i$ are known or estimated (Louvain, Leiden, spectral clustering):
\begin{equation}
    \tau_{ij} =
    \begin{cases}
        \tau_{\mathrm{in}}  & \text{if } c_i = c_j, \\
        \tau_{\mathrm{out}} & \text{if } c_i \neq c_j,
    \end{cases}
    \qquad \tau_{\mathrm{in}} \gg \tau_{\mathrm{out}}.
\end{equation}
Example: $\tau_{\mathrm{in}} \in [0.3, 0.5]$, $\tau_{\mathrm{out}} \in [0.01, 0.05]$.

\subsection{Structural Heuristics (No Labels)}

\begin{table}[h]
\centering
\caption{Topology-based signals for $\tau_{ij}$.}
\begin{tabular}{@{}lll@{}}
\toprule
\textbf{Signal} & \textbf{Rule of thumb} & \textbf{Rationale} \\
\midrule
Common neighbors / Jaccard & $\tau_{ij} \propto \mathrm{Jaccard}(i,j)$ & Shared neighborhood $\Rightarrow$ same community \\
Edge betweenness            & $\tau_{ij} \propto 1 - \mathrm{norm}(\mathrm{betweenness}_{ij})$ & Bridges get small $\tau$ \\
Effective resistance        & $\tau_{ij} \propto \exp(-R_{ij}/T)$ & Graphically distant pairs mix slowly \\
Conductance / cut edges     & Small $\tau$ on cross-cut edges         & Direct community-boundary control \\
\bottomrule
\end{tabular}
\end{table}

A practical template:
\begin{equation}
    \tau_{ij} = \tau_{\min} + (\tau_{\max}-\tau_{\min})\cdot \sigma\!\left(\frac{\mathrm{CN}_{ij}-\mu}{\sigma}\right),
\end{equation}
where $\mathrm{CN}_{ij}$ is the common-neighbor count.

\subsection{Feature Homophily}

Align with DGAD's feature-space view:
\begin{equation}
    \tau_{ij} = \tau_{\max}\cdot \exp\!\left(-\frac{\|x_i - x_j\|^2}{T}\right).
\end{equation}
Use initial features $x^{(0)}$ for a fixed $\tau_{ij}$, or current $x^{(k)}$ for dynamic $\tau$. Attention $\alpha_{ij}$ can govern \emph{who} to listen to; $\tau_{ij}$ governs \emph{how fast} to diffuse along each edge.

\subsection{Decoupling from Attention}

To avoid double amplification on already-strong edges:
\begin{equation}
    \tau_{ij} = \tau_{\mathrm{base}} \cdot f_{ij}, \quad f_{ij} \in [0,1],
\end{equation}
where $f_{ij}$ comes from structure or features, not directly from $\alpha_{ij}$. A per-node budget normalization is also possible:
\begin{equation}
    \tau_{ij} = \tau_{\max}\cdot \frac{f_{ij}}{\sum_{j'} \alpha_{ij'} f_{ij'}}\cdot \alpha_{ij}.
\end{equation}

\subsection{Learnable \texorpdfstring{$\tau_{ij}$}{tau\_ij}}

For end-to-end training:
\begin{equation}
    \tau_{ij} = \sigma\!\left(w^\top [h_i \| h_j \| e_{ij}] + b\right).
\end{equation}
Regularize large $\tau$ on predicted bridge edges, or distill from heuristic teachers $f_{ij}$.

\subsection{Step Schedule (Orthogonal but Complementary)}

Even with fixed $\tau_{ij}$, decay over steps slows late-stage global mixing:
\begin{equation}
    \tau_{ij}^{(k)} = \tau_{ij}\cdot \gamma^k, \quad \gamma < 1.
\end{equation}

\section{Recommended Progression}

\begin{enumerate}
    \item Fixed graph + Louvain communities: $\tau_{\mathrm{in}}=0.4$, $\tau_{\mathrm{out}}=0.02$.
    \item No community labels: $\tau_{ij} \propto \mathrm{Jaccard}(i,j)$ or $\exp(-\|x_i-x_j\|^2/T)$.
    \item Joint with attention: $\alpha$ for weights, $\tau$ for rates; cap inter-community $\tau$ aggressively.
    \item Supervised tasks: learnable $\tau$ with bridge regularization.
\end{enumerate}

Monitor:
\begin{itemize}[noitemsep]
    \item \textbf{Local signal}: rising cosine similarity \emph{within} communities;
    \item \textbf{Global signal}: \emph{between}-community similarity stays low or rises slowly.
\end{itemize}

\section{Comparison with Global \texorpdfstring{$\tau$}{tau}}

\begin{table}[h]
\centering
\caption{Global vs.\ per-edge step sizes.}
\begin{tabular}{@{}cccc@{}}
\toprule
\textbf{Method} & \textbf{Promote local} & \textbf{Slow global} & \textbf{Cost} \\
\midrule
Reduce $K$           & medium & strong & risk of under-smoothing \\
Reduce global $\tau$ & weak   & strong & slows intra-community mixing too \\
\textbf{Per-edge $\tau_{ij}$} & \textbf{strong} & \textbf{strong} & $O(|E|)$ storage/compute \\
\bottomrule
\end{tabular}
\end{table}

Per-edge $\tau$ allows large steps on intra-community edges and small steps on inter-community edges simultaneously---avoiding a single global $\tau$ that sacrifices one scale for the other.

%==============================================================================
\chapter{\texorpdfstring{$\tau_{ij}$}{tau\_ij} from KNN Graph Asymmetry}
%==============================================================================

\section{Asymmetry in Feature-Space KNN Graphs}

When building graphs via \texttt{knn\_graph}, each node $i$ selects $k$ nearest neighbors in feature space, creating directed edges $i \to j$. In \texttt{GNDLayer}, \texttt{edge\_index[0]} is the aggregation target, so node $i$ receives messages from its chosen neighbors.

Two degree quantities differ:

\begin{table}[h]
\centering
\caption{KNN degree asymmetry (aggregation convention).}
\begin{tabular}{@{}lll@{}}
\toprule
\textbf{Quantity} & \textbf{Meaning} & \textbf{On KNN graphs} \\
\midrule
In-degree (aggregation) $d^{\mathrm{in}}_i$ & Edges for which $i$ is the target & \textbf{Fixed $\approx k$} \\
Reverse degree $d^{\mathrm{rev}}_j$ & Count of nodes that select $j$ as a neighbor & \textbf{Varies with local density} \\
\bottomrule
\end{tabular}
\end{table}

High $d^{\mathrm{rev}}_j$ indicates that $j$ lies in a dense region of feature space (many points include $j$ in their KNN list). Low $d^{\mathrm{rev}}_j$ suggests peripheral, isolated, or boundary locations. This asymmetry is a free structural signal for designing $\tau_{ij}$ without community labels.

\section{Proposed Formulas}

For edge $(i \leftarrow j)$ (node $i$ receives from neighbor $j$):
\begin{equation}
    \tau_{ij} = \tau_{\min} + (\tau_{\max}-\tau_{\min})\cdot g\!\left(d^{\mathrm{rev}}_i,\, d^{\mathrm{rev}}_j\right).
\end{equation}

\subsection{Scheme 1: Source Density}

\begin{equation}
    g = \frac{d^{\mathrm{rev}}_j - d_{\min}}{d_{\max} - d_{\min}}.
\end{equation}
Messages from dense-region nodes diffuse faster.

\subsection{Scheme 2: Co-Density (Conservative)}

\begin{equation}
    g = \min\!\left(\frac{d^{\mathrm{rev}}_i}{d_{\max}},\, \frac{d^{\mathrm{rev}}_j}{d_{\max}}\right).
\end{equation}
Large $\tau$ only when \emph{both} endpoints are dense---boundary edges are automatically suppressed.

\subsection{Scheme 3: Density Similarity}

\begin{equation}
    g = \exp\!\left(-\frac{\bigl|d^{\mathrm{rev}}_i - d^{\mathrm{rev}}_j\bigr|^2}{T}\right).
\end{equation}
Similar reverse degrees (typical intra-community pairs) receive large $\tau$; mismatched pairs (dense core to sparse periphery) receive small $\tau$.

\subsection{Scheme 4: Mutual KNN + Co-Density (Recommended)}

Cross-community edges in asymmetric KNN are often \textbf{one-way}; intra-community edges are more often \textbf{mutual}:
\begin{equation}
    \tau_{ij} = \tau_{\max}\cdot \mathbb{1}[i \in \mathrm{KNN}(j)]\cdot \frac{\sqrt{d^{\mathrm{rev}}_i \cdot d^{\mathrm{rev}}_j}}{d_{\max}}.
\end{equation}

\begin{itemize}[noitemsep]
    \item mutual + high reverse degree $\Rightarrow$ large $\tau$: fast \textbf{local} equilibration inside dense communities;
    \item non-mutual or low density $\Rightarrow$ small $\tau$: slow \textbf{global} coupling.
\end{itemize}

\section{Integration with Diffusion}

Substitute $\tau_{ij}$ into Equation~\eqref{eq:edge-tau}. Precompute $d^{\mathrm{rev}}$ from \texttt{edge\_index} in $O(|E|)$ once per graph build. When \texttt{edge\_rewire=True}, recompute after each KNN reconstruction.

\section{Qualitative Picture}

\begin{quote}
\emph{Dense community interior}: high $d^{\mathrm{rev}}$ $\Rightarrow$ large $\tau$ $\Rightarrow$ fast local equilibrium.\\
\emph{Boundaries and bridges}: low or asymmetric $d^{\mathrm{rev}}$ $\Rightarrow$ small $\tau$ $\Rightarrow$ delayed global equilibrium.
\end{quote}

Per-edge diffusion rates follow the geometric role of each KNN edge rather than a single global step size.

\section{Caveats and Limitations}

\begin{enumerate}
    \item \textbf{Dense--dense bridges}: two community cores can both have high $d^{\mathrm{rev}}$. Schemes~1--2 alone may assign large $\tau$ to inter-community bridges; mutual-KNN (Scheme~4) or density-similarity (Scheme~3) mitigates this.
    \item \textbf{Dependence on $k$}: larger KNN $k$ shifts all $d^{\mathrm{rev}}$ upward; prefer quantile or $d_{\max}$ normalization over raw counts.
    \item \textbf{Dynamic rewiring}: $d^{\mathrm{rev}}$ changes when KNN is rebuilt each step; refresh $\tau_{ij}$ accordingly.
    \item \textbf{Fixed non-KNN graphs} (e.g., citation networks): reverse-degree asymmetry is absent; use Jaccard, community detection, or feature distance instead.
\end{enumerate}

\section{Practical Default}

A strong starting formula:
\begin{equation}
    \tau_{ij} = \tau_{\max}\cdot \mathbb{1}[i \in \mathrm{KNN}(j)]\cdot \exp\!\left(-\frac{\bigl|d^{\mathrm{rev}}_i - d^{\mathrm{rev}}_j\bigr|^2}{T}\right),
\end{equation}
with $\tau_{\max} \in [0.3, 0.5]$ and $T$ set to one or two times the variance of $\{d^{\mathrm{rev}}_i\}$. Success is indicated when intra-community embedding similarity rises quickly while inter-community similarity remains comparatively low.

%==============================================================================
\chapter{Summary}
%==============================================================================

DGAD diffusion produces \textbf{local} and \textbf{global equilibrium effects} at two complementary levels:
\begin{enumerate}
    \item \textbf{Hop scale}: one-hop neighbor mixing vs.\ multi-hop graph-wide harmonization;
    \item \textbf{Community scale}: fast intra-community consensus vs.\ slow inter-community coupling on structured graphs.
\end{enumerate}

For community separation, favor settings where local equilibria have formed inside communities but global equilibration across communities is still incomplete. Global knobs ($K$, $\tau$) provide coarse control; \textbf{per-edge $\tau_{ij}$} provides fine control---assigning large rates to homophilous, dense, mutual-KNN edges and small rates to sparse, asymmetric, or bridge-like edges.

On KNN graphs, fixed in-degree and variable reverse degree offer a particularly natural, label-free signal: \emph{popularity as a neighbor} proxies local feature density and can be turned directly into a heterogenous diffusion schedule that promotes local equilibrium while retarding global equilibrium.

\end{document}
